\documentclass[conference]{IEEEtran}
\IEEEoverridecommandlockouts

% -----------------------------------------------------------------------
% Phase II Final IEEE Report — MCA Project 7
% 1-MUX: Direct-Mapped vs. 2-Way Set-Associative Cache with LRU
% Team 5 — MCA Programme, Rishihood University, 2024--25
% -----------------------------------------------------------------------
% Build:  tectonic report/phase2/main.tex -o report/phase2/
% Status: DRAFT — FPGA synthesis results are PENDING (Kushal Sarkar)
% -----------------------------------------------------------------------

\usepackage{cite}
\usepackage{amsmath,amssymb,amsfonts}
\usepackage{algorithmic}
\usepackage{graphicx}
\usepackage{textcomp}
\usepackage{xcolor}
\usepackage{booktabs}
\usepackage{array}
\usepackage{url}
\usepackage{verbatim}

\def\BibTeX{{\rm B\kern-.05em{\sc i\kern-.025em b}\kern-.08em
    T\kern-.1667em\lower.7ex\hbox{E}\kern-.125emX}}

% -----------------------------------------------------------------------
\begin{document}

\title{RTL Design and Verification of a 1-Mux Direct-Mapped and
       2-Way Set-Associative Cache with LRU Replacement:\\
       Phase II Final Report}

\author{%
  \IEEEauthorblockN{Krishna Pramod Dubale}
  \IEEEauthorblockA{\textit{MCA Programme}\\
    Rishihood University, Sonipat, India\\
    krishna.dubale2024@nst.rishihood.edu.in}
  \and
  \IEEEauthorblockN{Apoorva Choudhary}
  \IEEEauthorblockA{\textit{MCA Programme}\\
    Rishihood University, Sonipat, India\\
    apoorva.choudhary2024@nst.rishihood.edu.in}
  \and
  \IEEEauthorblockN{Prachee Dhar}
  \IEEEauthorblockA{\textit{MCA Programme}\\
    Rishihood University, Sonipat, India\\
    prachee.dhar2024@nst.rishihood.edu.in}
  \and
  \IEEEauthorblockN{Kushal Sarkar}
  \IEEEauthorblockA{\textit{MCA Programme}\\
    Rishihood University, Sonipat, India\\
    kushal.sarkar2024@nst.rishihood.edu.in}
  \and
  \IEEEauthorblockN{Vani Rudra}
  \IEEEauthorblockA{\textit{MCA Programme}\\
    Rishihood University, Sonipat, India\\
    vani.rudra2024@nst.rishihood.edu.in}
}

\maketitle

% -----------------------------------------------------------------------
\begin{abstract}
This report presents the complete Phase II evaluation of a Verilog RTL
implementation comparing a direct-mapped cache baseline and a 2-way
set-associative cache with true 1-bit LRU replacement.
Both designs target equal 128-byte capacity, 16-byte blocks, and an
identical synchronous handshake interface ({\tt req\_valid}/{\tt req\_ready}
master, {\tt resp\_valid}/{\tt resp\_hit} slave).
Phase I established the RTL, interface, and architecture; Phase II extends
the work with a four-workload comparative benchmark suite (89 total trace
requests across conflict, spatial-locality, extended-conflict, and
sequential traces), formal verification suites totalling 618 assertions,
latency characterisation under a combinational lower-memory model, and
modelled Average Memory Access Time (AMAT) analysis.

Key findings are: (1)~2-way LRU achieves up to 87.5\% hit rate on
extended two-block conflict workloads where direct-mapped degrades to 0\%;
(2)~both designs achieve identical hit rates on spatial-locality and
sequential workloads; (3)~a 3-block repeating workload exposes a set-index
anomaly in which the direct-mapped design outperforms the 2-way cache due
to its eight-set geometry; (4)~simulation-measured acceptance-to-response
latency averages 2.506 cycles per request for the 2-way cache versus
2.584 cycles for direct-mapped.
FPGA synthesis results targeting Xilinx/AMD hardware are \textbf{PENDING}
and will be integrated when available.
\end{abstract}

\begin{IEEEkeywords}
cache memory, set-associative, LRU replacement, RTL design, Verilog,
verification, AMAT, FPGA synthesis, direct-mapped, conflict miss
\end{IEEEkeywords}

% -----------------------------------------------------------------------
\section{Introduction}
\label{sec:intro}

Cache memory sits at the intersection of processor speed and memory
bandwidth, exploiting temporal and spatial locality to reduce average
memory access time \cite{smith1982}.
Choosing between direct-mapped and set-associative organisations requires
careful evaluation of workload sensitivity, silicon area, and timing
characteristics \cite{hill1988,przybylski1988}.

This project designs two cache controllers in synthesisable Verilog and
subjects them to a systematic verification and benchmarking regime.
Phase I \cite{ieee-verilog2001} produced the RTL modules, testbenches,
interface specification, and an initial five-access conflict trace
demonstrating the associativity advantage.
Phase II deepens the evaluation with four workload traces covering
pathological conflict, spatial locality, and capacity behaviour, together
with 618-assertion correctness suites and latency modelling.

The remainder of this report is structured as follows:
Section~\ref{sec:litreview} surveys related work;
Section~\ref{sec:solution} describes the RTL architecture and design
decisions;
Section~\ref{sec:results} presents functional verification, benchmark, and
latency results;
Section~\ref{sec:contributions} records individual contributions;
Section~\ref{sec:conclusion} concludes with directions for future work.

% -----------------------------------------------------------------------
\section{Literature Review}
\label{sec:litreview}

Smith's landmark survey \cite{smith1982} characterised cache behaviour
across working-set size, block size, and associativity dimensions, showing
that set-associative designs consistently reduce conflict misses when the
working set exceeds the number of direct-mapped sets.
Hill \cite{hill1988} argued the counter-case: for identical silicon area,
direct-mapped caches achieve a faster cycle time, which may compensate for
their higher miss rate.

Przybylski, Horowitz, and Hennessy \cite{przybylski1988} formalised the
performance trade-off in terms of AMAT:
\begin{equation}
  \text{AMAT} = t_{\text{hit}} + (1 - h) \times M_p
  \label{eq:amat}
\end{equation}
where $t_{\text{hit}}$ is the hit-path latency, $h$ is the hit rate, and
$M_p$ is the miss penalty.
Our Phase II evaluation uses this model with $t_{\text{hit}} = 2$ cycles
and $M_p = 20$ cycles as a representative DRAM refill penalty.

Jouppi \cite{jouppi1990} demonstrated that victim buffers and prefetch
streams can reclaim much of the conflict-miss cost in direct-mapped
designs; our implementation does not include these but provides a clean
baseline for such extensions.
Al-Zoubi et al.\ \cite{al-zoubi2004} evaluated replacement policies on
SPEC CPU2000, confirming LRU's strength on workloads with high temporal
locality and noting its limitation on cyclic traces that exceed associativity.
This matches our Trace E finding (3-block cycle exposing the 2-way LRU
limitation).

Agarwal, Hennessy, and Horowitz \cite{agarwal1989} studied cache behaviour
under multiprogramming, highlighting that set-index aliasing is a primary
source of conflict misses in OS workloads — directly motivating our
extended Trace C (16-access alternating conflict) which demonstrates 87.5\%
versus 0\% hit rates.

FPGA implementations of cache memories commonly target Xilinx Block RAM
(BRAM) primitives for the data array and distributed LUT-RAM for small
tag arrays, following guidelines from the Vivado synthesis user guide
\cite{vivado2023}.
Our RTL uses Verilog register arrays, which Vivado infers as distributed
RAM or BRAM depending on depth and read-port count; exact utilisation
figures are \textbf{PENDING} from Phase II synthesis experiments.

% -----------------------------------------------------------------------
\section{Proposed Solution}
\label{sec:solution}

\subsection{Design Goals}
\label{subsec:goals}

The design targets four goals: (1)~a read-only, synthesisable RTL
implementation suitable for FPGA deployment; (2)~an equal-capacity
baseline comparison that isolates the associativity variable;
(3)~a fully-verified handshake interface shared between both modules; and
(4)~a deterministic LRU policy implementable with a single bit per set.

\subsection{Module Hierarchy}
\label{subsec:hierarchy}

Three RTL files implement the design:

\begin{itemize}
  \item {\tt cache\_2way.v} — top-level 2-way set-associative controller;
        instantiates {\tt cache\_2way\_store}.
  \item {\tt cache\_2way\_store.v} — parameterised storage array holding
        tag, valid, data, and 1-bit LRU state for four sets.
  \item {\tt cache\_direct.v} — self-contained direct-mapped controller
        with eight sets.
\end{itemize}

Both controllers expose identical ports: {\tt clk}, {\tt rst},
{\tt req\_addr[31:0]}, {\tt req\_valid}, {\tt req\_ready},
{\tt resp\_data[31:0]}, {\tt resp\_hit}, {\tt resp\_valid}.

\subsection{Address Decomposition}
\label{subsec:addr}

Both caches store 128~bytes in 8 blocks of 16~bytes each.
The 32-bit byte address is partitioned differently depending on associativity,
as shown in Table~\ref{tab:addr}.

\begin{table}[h]
  \caption{Address Field Decomposition}
  \label{tab:addr}
  \centering
  \small
  \begin{tabular}{lcccc}
    \toprule
    Design & Tag & Set Index & Word & Byte Offset \\
    \midrule
    2-Way  & {[}31:6{]} 26b & {[}5:4{]} 2b & {[}3:2{]} 2b & {[}1:0{]} 2b \\
    Direct & {[}31:7{]} 25b & {[}6:4{]} 3b & {[}3:2{]} 2b & {[}1:0{]} 2b \\
    \bottomrule
  \end{tabular}
\end{table}

The key geometric consequence is that address bit~6 is a \emph{set-index
bit} in direct-mapped but a \emph{tag bit} in 2-way, which causes
{\tt 0x00} and {\tt 0x40} to reside in different sets under direct-mapped
but to conflict in Set~0 under 2-way. This distinction drives the Trace~E
set-index anomaly result (\S\ref{subsec:anomaly}).

\subsection{Controller FSM}
\label{subsec:fsm}

Both controllers implement an identical 4-state Mealy FSM:

\begin{verbatim}
  IDLE  --req_valid--> LOOKUP
  LOOKUP --hit------> RESP
  LOOKUP --miss-----> REFILL
  REFILL ------------> RESP
  RESP  ------------> IDLE
\end{verbatim}

\noindent
In {\tt STATE\_IDLE} the controller asserts {\tt req\_ready = 1} and
latches {\tt req\_addr} into {\tt saved\_addr} on acceptance.
{\tt STATE\_LOOKUP} performs combinational tag comparison against both
ways (2-way) or the single set (direct).
On a miss, {\tt STATE\_REFILL} writes the lower-memory block into the
appropriate way, updating the LRU bit.
{\tt STATE\_RESP} asserts {\tt resp\_valid}, {\tt resp\_hit}, and
{\tt resp\_data}, then returns to {\tt IDLE}.

\subsection{LRU Replacement Mechanics}
\label{subsec:lru}

The 2-way cache maintains a single LRU bit per set stored in
{\tt cache\_2way\_store}.
On a \textbf{hit} in Way~$w$: {\tt lru\_bit $\leftarrow$ 1 - $w$}
(point to the \emph{other} way as next victim).
On a \textbf{miss}: the way indicated by {\tt lru\_bit} is overwritten,
then {\tt lru\_bit} is flipped.

\textbf{Worked example} (addresses {\tt 0x00}, {\tt 0x80}, {\tt 0x100},
all mapping to Set~0):
\begin{enumerate}
  \item {\tt LOAD 0x00} — cold miss, fill Way~0; lru $\leftarrow$ 1 (evict Way~1 next).
  \item {\tt LOAD 0x80} — cold miss, fill Way~1; lru $\leftarrow$ 0 (evict Way~0 next).
  \item {\tt LOAD 0x00} — \emph{hit} in Way~0; lru $\leftarrow$ 1.
  \item {\tt LOAD 0x100} — miss, evict Way~1 (lru=1 $\rightarrow$ evict Way~1);
        lru $\leftarrow$ 0.
  \item {\tt LOAD 0x00} — \emph{hit} in Way~0 (Way~0 was not evicted). lru $\leftarrow$ 1.
\end{enumerate}
The 2-way cache retains {\tt 0x00} across all five accesses (2 hits out
of 5), whereas the direct-mapped cache achieves 0 hits on the same trace.

\subsection{Read-Only Transaction Model}
\label{subsec:readonly}

Phase I and Phase II implement a \emph{read-only} cache: no write-back,
write-allocate, or dirty bits exist in the RTL.
The combinational lower-memory model ({\tt u\_mem} in testbenches) returns
data in the same clock cycle as the refill request, eliminating DRAM
latency from simulation measurements.
This is a deliberate simplification: it allows the FSM and hit/miss logic
to be validated independently of memory timing, and it means measured
simulation cycles represent \emph{hardware pipeline depth} rather than
wall-clock performance.
Section~\ref{subsec:amat} discusses how AMAT modelling restores the
realistic penalty.

% -----------------------------------------------------------------------
\section{Results and Discussion}
\label{sec:results}

\subsection{Functional Verification}
\label{subsec:verification}

Table~\ref{tab:suites} summarises the five verification suites run across
the Phase~I and Phase~II lifecycle.
All 618 assertions pass with Icarus Verilog 12.0 \cite{icarus-verilog}
on the {\tt main} branch.

\begin{table}[h]
  \caption{Verification Suite Summary (All Passing)}
  \label{tab:suites}
  \centering
  \small
  \begin{tabular}{llr}
    \toprule
    Suite & Owner & Assertions \\
    \midrule
    LRU unit tests       & Krishna Pramod Dubale   &  51 \\
    Controller FSM tests & Apoorva Choudhary       & 141 \\
    Cache comparison     & Prachee Dhar            &  52 \\
    System regression    & Kushal Sarkar           & 344 \\
    Waveform demo        & Kushal Sarkar           &  30 \\
    \midrule
    \textbf{Total}       &                         & \textbf{618} \\
    \bottomrule
  \end{tabular}
\end{table}

\subsection{Benchmark Workload Results}
\label{subsec:benchmarks}

Four trace workloads were evaluated by Prachee Dhar using
{\tt tb/tb\_cache\_compare.sv} against on-disk traces in
{\tt tb/traces/}.
Table~\ref{tab:hitrates} reports hit rates and Table~\ref{tab:amat_table}
reports modelled AMAT.

\begin{table}[h]
  \caption{Hit Rate Comparison by Workload Trace}
  \label{tab:hitrates}
  \centering
  \small
  \begin{tabular}{lp{3.0cm}cc}
    \toprule
    Trace & Description & 2-Way Hit\% & Direct Hit\% \\
    \midrule
    A & 5-access conflict (0x00/0x80 thrash)   & 40.0  & 0.0  \\
    B & 6-access within-block spatial locality  & 66.7  & 66.7 \\
    C & 16-access extended two-block conflict   & 87.5  & 0.0  \\
    D & 32-access full-capacity sequential      & 75.0  & 75.0 \\
    E & 10-access 3-block cycle (0x00/0x40/0x80)& 0.0   & 30.0 \\
    \bottomrule
  \end{tabular}
\end{table}

\begin{table}[h]
  \caption{Modelled AMAT ($t_{\text{hit}}=2$ cycles, $M_p=20$ cycles)}
  \label{tab:amat_table}
  \centering
  \small
  \begin{tabular}{lcc}
    \toprule
    Trace & 2-Way AMAT (cycles) & Direct AMAT (cycles) \\
    \midrule
    A &  14.00 & 22.00 \\
    B &   8.67 &  8.67 \\
    C &   4.50 & 22.00 \\
    D &   7.00 &  7.00 \\
    E &  22.00 & 16.00 \\
    \bottomrule
  \end{tabular}
\end{table}

\subsection{Conflict and Spatial-Locality Findings}
\label{subsec:conflict}

\textbf{2-way wins on conflict workloads.}
Traces A and C expose the fundamental weakness of direct-mapped placement:
any two distinct addresses mapping to the same set cause mutual eviction on
every alternating access.
The 2-way cache eliminates this for working sets of up to two blocks per
set.
On the 16-access Trace~C, the 2-way cache achieves 87.5\% (14/16) versus
0\% (0/16) — an AMAT advantage of 17.5 cycles per access under the 20-cycle
penalty model.

\textbf{Both tie on spatial-locality and sequential workloads.}
Traces B and D show identical hit rates across both designs.
The refill granularity is one 16-byte block regardless of associativity,
so consecutive word accesses within a block hit identically.
This confirms that associativity provides no additional benefit when the
access pattern fits naturally within a single block.

\subsection{Set-Index Anomaly}
\label{subsec:anomaly}

Trace~E (addresses {\tt 0x00}, {\tt 0x40}, {\tt 0x80} cycling) reveals
that direct-mapped \emph{outperforms} 2-way (30\% vs.\ 0\%) on 3-block
repeating cycles.
In direct-mapped, bit~6 of the address is a set-index bit:
{\tt 0x40} maps to Set~4 while {\tt 0x00} and {\tt 0x80} map to Set~0,
so only the two-block conflict in Set~0 thrashes, and {\tt 0x40} hits
stably.
In 2-way, all three addresses map to Set~0 (bit~6 is a tag bit), which
overflows the two-way capacity and triggers continuous LRU eviction of
all three blocks.

This result underscores Hill's argument \cite{hill1988}: the relationship
between associativity and hit rate is workload-dependent, not monotone.

\subsection{Latency Characterisation}
\label{subsec:latency}

Under the combinational memory model, hit latency is 2 clock cycles
(IDLE $\rightarrow$ LOOKUP $\rightarrow$ RESP) and miss latency is 3 cycles
(LOOKUP $\rightarrow$ REFILL $\rightarrow$ RESP).
Across 89 total trace requests, the weighted average acceptance-to-response
latency is:
\begin{itemize}
  \item \textbf{2-way cache:} 2.506 cycles/request.
  \item \textbf{Direct-mapped:} 2.584 cycles/request.
\end{itemize}

Modelled AMAT restores the 20-cycle DRAM penalty as described in
\S\ref{subsec:amat}.
On conflict-heavy workloads, the 2-way AMAT is up to 17.5 cycles lower
per access, far outweighing the negligible 0.078-cycle raw pipeline
difference.

\subsection{AMAT Analysis}
\label{subsec:amat}

Using Equation~\ref{eq:amat} with $t_{\text{hit}} = 2.0$ cycles and
$M_p = 20$ cycles:
\begin{align*}
  \text{AMAT}_{\text{2way,C}}   &= 2.0 + (1 - 0.875) \times 20 = 4.5 \text{ cycles} \\
  \text{AMAT}_{\text{direct,C}} &= 2.0 + (1 - 0.000) \times 20 = 22.0 \text{ cycles}
\end{align*}
The 2-way cache provides a $\approx\!4.9\times$ AMAT advantage on
extended two-block conflict (Trace~C), dropping to $1.0\times$ on
spatial-locality traces.

\subsection{Limitations}
\label{subsec:limits}

\begin{itemize}
  \item \textbf{Combinational memory model:} The testbench memory returns
        data in zero simulation cycles, eliminating DRAM latency.
        AMAT figures are therefore modelled rather than measured.
  \item \textbf{Read-only operation:} No dirty bits, write-back, or
        write-allocate paths are implemented.
        Real workloads with stores would require additional FSM states.
  \item \textbf{Single outstanding request:} The FSM blocks in
        STATE\_REFILL, so pipelining or out-of-order requests are
        not supported.
  \item \textbf{FPGA timing:} Measured simulation cycles cannot be
        converted to wall-clock time without synthesis-reported $F_{\max}$.
        FPGA results are \textbf{PENDING} (\S\ref{subsec:fpga}).
\end{itemize}

\subsection{FPGA Synthesis Results}
\label{subsec:fpga}

\textbf{[PENDING — Kushal Sarkar to provide synthesis artifacts.]}

Target device: Xilinx/AMD Artix-7 (or equivalent), Vivado 2023.2,
clock constraint 100~MHz ($T_{\text{clk}} = 10$~ns)
\cite{vivado2023}.
Table~\ref{tab:fpga} will be completed with verified synthesis output.

\begin{table}[h]
  \caption{FPGA Synthesis Utilisation and Timing (PENDING)}
  \label{tab:fpga}
  \centering
  \small
  \begin{tabular}{lccc}
    \toprule
    Metric & 2-Way & Direct & Unit \\
    \midrule
    LUT count            & --- & --- & LUT4 \\
    Flip-flop count      & --- & --- & FF   \\
    BRAM usage           & --- & --- & 36Kb tiles \\
    $W_{\text{NS}}$ (slack) & --- & --- & ns \\
    $F_{\max}$ (reported) & --- & --- & MHz \\
    \bottomrule
  \end{tabular}
\end{table}

% -----------------------------------------------------------------------
\section{Individual Contributions}
\label{sec:contributions}

\begin{itemize}
  \item \textbf{Krishna Pramod Dubale} — Authored {\tt cache\_2way\_store.v}
        (tag, valid, data, LRU arrays) and {\tt tb/tb\_lru\_unit.sv} (51
        assertions covering LRU update on hit, miss, and double-eviction
        edge cases). Phase~II branch: {\tt phase2/krishna-lru-verification}.

  \item \textbf{Apoorva Choudhary} — Authored {\tt tb/tb\_cache\_2way.sv}
        (141 assertions covering all four FSM state transitions, hit/miss
        flag correctness, {\tt resp\_data} correctness on spatial accesses,
        and {\tt req\_ready} deassertion during refill).
        Phase~II branch: {\tt phase2/apoorva-controller-verification}.

  \item \textbf{Prachee Dhar} — Designed and executed the four-workload
        comparative benchmark suite ({\tt tb/tb\_cache\_compare.sv},
        on-disk traces in {\tt tb/traces/}, {\tt scripts/run\_cache\_compare.sh}).
        Documented results in {\tt docs/phase2/cache-benchmarks.md}.
        Phase~II branch: {\tt phase2/prachee-cache-benchmarks}.

  \item \textbf{Kushal Sarkar} — System-level integration: 344 regression
        assertions ({\tt tb/tb\_system\_regression.sv}), waveform
        evidence ({\tt docs/evidence/cache\_2way\_demo.vcd}),
        HANDOFF.md entries, and ongoing FPGA synthesis.
        Phase~II synthesis results PENDING.

  \item \textbf{Vani Rudra} — Phase~II final IEEE report (this document),
        architecture and FSM figures, address decomposition tables,
        AMAT analysis and results discussion, and interview handout
        ({\tt docs/phase2/interview.md}).
        Phase~II branch: {\tt phase2/vani-final-report}.
\end{itemize}

% -----------------------------------------------------------------------
\section{Conclusion}
\label{sec:conclusion}

This project has demonstrated, through 618 passing assertions and four
workload benchmarks, that 2-way set-associative caching with 1-bit LRU
replacement eliminates conflict thrashing at a modest area cost relative
to a direct-mapped design of equal capacity.
The key quantitative result is that the 2-way cache achieves up to 87.5\%
hit rate on extended two-block conflict workloads, while direct-mapped
degrades to 0\%  — an AMAT improvement of $\approx\!4.9\times$ under a
realistic 20-cycle DRAM penalty.

The set-index anomaly on the 3-block cycle trace is an important
qualification: associativity is beneficial only when the working set fits
within the per-set capacity.
Workloads with purely spatial locality benefit neither design over the other;
the block-level refill granularity equalises hit rates in that regime.

Future work should address write operations (write-back or write-through),
multi-level cache hierarchies, out-of-order request handling, and FPGA
physical implementation with accurate timing closure.
Integrating Kushal's FPGA synthesis data will convert the modelled AMAT
figures into wall-clock performance comparisons and confirm whether the
combinational-memory FSM meets 100~MHz timing.

% -----------------------------------------------------------------------
\section*{Acknowledgment}

The authors thank the faculty of the MCA Programme, Rishihood University,
for project guidance, and acknowledge the Icarus Verilog open-source
community \cite{icarus-verilog} for the simulation infrastructure used
throughout this project.

% -----------------------------------------------------------------------
\bibliographystyle{IEEEtran}
\bibliography{refs}

\end{document}
